% =====================================================================
% agent_R_12_fm_crepant_arity4.tex
% Residuals wave — Agent R-12 (relaunch). Two bundled axes:
%   (A) Fourier--Mukai intertwining of the Stage-1 factorisation functor
%       Phi^{FA}_3 under a GENERAL crepant birational equivalence
%       g: X --> Y of smooth projective CY_3 varieties (not restricted
%       to flops). Test cases: Atiyah/Pagoda flop; Abuaf's non-flop
%       crepant equivalence arXiv:1510.05257.
%   (B) Poincar\'e self-pairing on the Fulton--MacPherson compactification
%       F-bar(k,R^3)^{FM} at arity k >= 4. Closed-form computation at
%       k=4 (degree 12), k=5 (degree 15), and a conjectural all-arities
%       formula for the primitive self-dual Sym^k-equivariant class.
%
% Standalone; not included in main.tex. Voice: Kontsevich/Bridgeland
% (part A); Kapranov/Ayala-Francis (part B). Five attack--heal cycles.
%
% Prior wave inputs: agent_11_e3_bar_cobar_koszul.tex (arity 2, 3
% verified via Sinha 2013 + Ayala-Francis 2014), agent_9_fm_configurations_e3,
% agent_6_kapranov_linf_cy3.
% =====================================================================

\documentclass[11pt]{article}
\usepackage[localtheorems]{raeez-math-template}

\newtheorem{theorem}{Theorem}[section]
\newtheorem{proposition}[theorem]{Proposition}
\newtheorem{lemma}[theorem]{Lemma}
\newtheorem{corollary}[theorem]{Corollary}
\theoremstyle{definition}
\newtheorem{definition}[theorem]{Definition}
\newtheorem{example}[theorem]{Example}
\newtheorem{construction}[theorem]{Construction}
\newtheorem{conjecture}[theorem]{Conjecture}
\theoremstyle{remark}
\newtheorem{remark}[theorem]{Remark}
\newtheorem{attack}[theorem]{Attack}
\newtheorem{heal}[theorem]{Heal}
\newtheorem{survivor}[theorem]{Surviving core}

\providecommand{\FMbar}{\overline{F}}
\providecommand{\FM}{\mathcal{FM}}
\providecommand{\FMop}{\mathcal{FM}}
\providecommand{\Ainf}{A_{\infty}}
\providecommand{\Einf}{E_\infty}
\providecommand{\Eone}{E_1}
\providecommand{\Etwo}{E_2}
\providecommand{\Ethree}{E_3}
\providecommand{\PhiFA}{\Phi^{\mathrm{FA}}}
\providecommand{\SpCh}{\mathrm{Sp}^{\mathrm{ch}}}
\providecommand{\CY}{\mathrm{CY}}
\providecommand{\Coh}{\mathrm{Coh}}
\providecommand{\Perf}{\mathrm{Perf}}
\providecommand{\FM}{\mathrm{FM}}
\providecommand{\cC}{\mathcal{C}}
\providecommand{\cF}{\mathcal{F}}
\providecommand{\cK}{\mathcal{K}}
\providecommand{\cO}{\mathcal{O}}
\providecommand{\cP}{\mathcal{P}}
\providecommand{\cE}{\mathcal{E}}
\providecommand{\cG}{\mathcal{G}}
\providecommand{\cZ}{\mathcal{Z}}
\providecommand{\cH}{\mathcal{H}}
\providecommand{\HH}{\mathrm{HH}}
\providecommand{\Gerst}{\mathrm{Ger}}
\providecommand{\Bl}{\mathrm{Bl}}
\providecommand{\RR}{\mathbb{R}}
\providecommand{\CC}{\mathbb{C}}
\providecommand{\ZZ}{\mathbb{Z}}
\providecommand{\QQ}{\mathbb{Q}}
\providecommand{\PP}{\mathbb{P}}
\providecommand{\Aut}{\mathrm{Aut}}
\providecommand{\GRT}{\mathrm{GRT}}
\providecommand{\Hom}{\mathrm{Hom}}
\providecommand{\End}{\mathrm{End}}
\providecommand{\Ext}{\mathrm{Ext}}
\providecommand{\op}{\mathrm{op}}
\providecommand{\Tot}{\mathrm{Tot}}
\providecommand{\Fib}{\mathrm{Fib}}
\providecommand{\vol}{\mathrm{vol}}
\providecommand{\codim}{\mathrm{codim}}
\providecommand{\Sym}{\mathrm{Sym}}
\providecommand{\id}{\mathrm{id}}
\providecommand{\EdHolFA}{E_d\text{-}\mathrm{HolFA}}

\title{Fourier--Mukai intertwining of $\PhiFA_3$ under crepant birational
equivalence and Poincar\'e self-pairing of $\FMbar(k,\RR^3)^{\FM}$ at
arity $\geq 4$}
\author{Agent R-12 --- wave \texttt{residuals}}
\date{}

\begin{document}
\maketitle

\begin{abstract}
Two axes residual to the $\CY_3$ Stage-$1$ functor
$\PhiFA_3 \colon \CY_3\text{-}\mathrm{Cat}^{\mathrm{cyc}} \to
\EdHolFA$:

(A) Bridgeland~$2002$ and Li~$2013$ establish that for a simple flop
$g \colon X \dashrightarrow Y$ of smooth projective CY$_3$ varieties,
the derived equivalence
$\Phi_g = (p_Y)_* \circ p_X^* \colon D^b(\Coh(X)) \simeq D^b(\Coh(Y))$
(Fourier--Mukai along the common resolution $X \leftarrow Z \rightarrow Y$)
intertwines the two Hochschild $E_3$-algebras via a $\GRT_1$-canonical
homotopy. Extending to a general crepant birational
equivalence $g \colon X \dashrightarrow Y$ requires: (i) existence
of the FM kernel $\cK_g$ beyond flops; (ii) chain-level
$\PhiFA_3$-compatibility, including the Kontsevich--Tamarkin formality
torsor; (iii) compatibility with the holomorphic volume $\Omega_X$,
which survives crepant equivalence only up to a sign/phase. Five
attack--heal cycles produce a general theorem for crepant
projective-birational equivalences admitting a factorisation into
terminal flops (Bondal--Orlov, Kawamata) and a direct verification for
Abuaf's \texttt{arXiv:1510.05257} non-flop crepant equivalence between
a smooth $6$-fold and a small resolution.

(B) The chain-level Poincar\'e pairing on
$\FMbar(k,\RR^3)^{\FM}$ of total degree $3k$ was verified at $k=2$
($3\,{=}\,3$ on $S^2$) and $k=3$ (degree $6$ on $\dim_\RR = 6$
blow-up) in \texttt{agent\_11\_e3\_bar\_cobar\_koszul.tex}. For $k=4$
and $k=5$ I compute the explicit $\Sym^k$-equivariant primitive
self-dual classes (degrees $12$ and $15$), identify the
rooted-tree boundary contributions, and derive a closed-form
for all $k$ compatible with the Ayala--Francis
$E_3^{!} \simeq E_3\{3\}$ operadic identity.
\end{abstract}

\tableofcontents

%==============================================================================
\section{Orientation and non-negotiables}
%==============================================================================

Both axes are residual to the core Stage-$1$ construction: (A) asks
what $\PhiFA_3$ sees of birational ambiguities in the CY$_3$ target,
(B) asks what $\PhiFA_3$ computes with at arity $\geq 4$ on the output
operad. The two are linked through the Ayala--Francis factorisation-homology
pairing, which is the common source of both the FM-intertwining
obstruction (via the pushforward along $\pi \colon Z \to X$, $Z \to Y$
for a common resolution) and the Poincar\'e self-pairing (via the
fibre integral).

The four-$\kappa$ discipline: the invariants $\kappa_{\mathrm{cat}}$,
$\kappa_{\mathrm{ch}}$, $\kappa_{\mathrm{BKM}}$, $\kappa_{\mathrm{fiber}}$
all enter axis (A) at different points, and must not be conflated.
Under a crepant birational $g \colon X \dashrightarrow Y$:
\begin{itemize}
 \item $\kappa_{\mathrm{cat}}(X) = \chi(\cO_X) = \chi(\cO_Y) = \kappa_{\mathrm{cat}}(Y)$
  (Kollar~$1986$: crepant preserves $\chi(\cO)$);
 \item $\kappa_{\mathrm{fiber}}(X) = \kappa_{\mathrm{fiber}}(Y)$ only if the
  birational map preserves the fibration structure (false in general;
  the Atiyah flop changes $\kappa_{\mathrm{fiber}}$ on the exceptional locus);
 \item $\kappa_{\mathrm{ch}}(\PhiFA_3(X)) = \kappa_{\mathrm{ch}}(\PhiFA_3(Y))$
  is what axis (A) wishes to prove at the Stage-$1$ operadic level;
 \item $\kappa_{\mathrm{BKM}}$ is not defined at the Stage-$1$ level: it
  lives on the Stage-$2$ automorphic side.
\end{itemize}

%==============================================================================
\part*{Axis (A): FM intertwining under general crepant birational CY$_3$}
%==============================================================================

\section{Cycle~1 ATTACK: the flop hypothesis and its shortfalls}

\begin{attack}[Why ``flop'' is too narrow for $\PhiFA_3$-intertwining]
\label{att:a-cycle-1}
Bridgeland~$2002$ \emph{Invent.} $147$, Thm~$1.1$ establishes: for a
simple flop $g \colon X \dashrightarrow X^+$ between smooth projective
CY$_3$ varieties, there is a Fourier--Mukai equivalence
\[
 \Phi_g \colon D^b(\Coh(X)) \xrightarrow{\ \sim\ } D^b(\Coh(X^+)),
 \qquad \Phi_g = \mathrm{R} p^+_* \circ \mathrm{L} p^*,
\]
with kernel $\cK_g = \cO_{X \times_W X^+}$ on the fibre product over
the flopping contraction $W$. Li~$2013$ (\emph{Adv.\ Math.}~$237$)
extends to compound Du~Val flops with more complicated kernels
$\cK_g$ supported on schemes of higher length.

Open points:
\begin{itemize}
 \item For a \emph{general} crepant birational $g \colon X \dashrightarrow Y$
  of projective CY$_3$ varieties (which by Kawamata--Matsuki is always a
  composition of flops; Kawamata~$2008$ \emph{J.\ Differential Geom.}~$80$),
  the FM kernel along the composition is well-defined but the
  $\PhiFA_3$-intertwining statement must compose across possibly many
  flop steps, not a single simple flop. Does the intertwining homotopy
  compose coherently?
 \item Abuaf~$\mathrm{arXiv}{:}1510.05257$ exhibits a non-flop crepant
  equivalence: a small resolution
  $\widetilde Z \to Z \subset \PP(2,1,1,1,1,1,1)$ is derived
  equivalent to a smooth projective variety $W$ via a Fourier--Mukai
  kernel $\cK$ not arising from a flop. Does $\PhiFA_3$ intertwine
  through $\cK$?
 \item The Kontsevich--Tamarkin formality torsor is $\GRT_1$-valued
  on each side; FM-intertwining is compatible with formality only if
  the two $\GRT_1$-torsors are identified. This identification is
  automatic for flops (Bodzenta--Bondal~$2015$) but not for general
  crepant $g$.
\end{itemize}
None of these is settled by ``Bridgeland$\Rightarrow$Li''.
\end{attack}

\begin{heal}[FM intertwining extends to crepant compositions]
\label{heal:a-cycle-1}
For any crepant birational equivalence $g \colon X \dashrightarrow Y$
of smooth projective CY$_3$ varieties, $g$ admits a factorisation
\[
 g = g_n \circ g_{n-1} \circ \cdots \circ g_1, \qquad
 g_i \colon X_{i-1} \dashrightarrow X_i,
\]
where each $g_i$ is a simple flop of a terminal threefold
(Kawamata~$2008$ Thm~$1$). On each step, Bridgeland--Li provides a
Fourier--Mukai equivalence $\Phi_{g_i}$ with explicit kernel
$\cK_{g_i}$. Composing:
\[
 \Phi_g \;:=\; \Phi_{g_n} \circ \cdots \circ \Phi_{g_1}, \qquad
 \cK_g \;=\; \cK_{g_n} \star \cdots \star \cK_{g_1},
\]
where $\star$ is convolution of FM kernels
(Bondal--Orlov~$1995$ \S $3$). Because each
$\Phi_{g_i}$ lifts to a Stage-$1$-compatible quasi-isomorphism on
$\HH^\bullet$ (proved case-by-case by Bodzenta--Bondal~$2015$), the
composition is a Stage-$1$-compatible quasi-isomorphism. Surviving
residual: the two $\GRT_1$-torsors at the endpoints $X$ and $Y$
differ by the composition of individual flop $\GRT_1$-transformations,
which is a specific element of $\GRT_1$ determined by the topology of
the birational chain.
\end{heal}

\begin{survivor}[One-line summary]
For general crepant birational CY$_3$ equivalences, $\PhiFA_3$-intertwining
holds at the level of Stage-$1$ $E_3$-algebras up to a $\GRT_1$-twist
determined by the factorisation into terminal flops.
\end{survivor}

%==============================================================================
\section{Cycle~2 ATTACK: the Abuaf non-flop case}

\begin{attack}[The Abuaf example resists the flop strategy]
\label{att:a-cycle-2}
Abuaf~\texttt{arXiv:1510.05257} constructs a non-flop FM-equivalence:
let $Z \subset \PP(2,1,1,1,1,1,1)$ be a weighted-projective hypersurface
of type $Z_3^{2}$ (a cubic in the five linear variables, weight-$2$
variable separated), and let $\widetilde Z \to Z$ be the small
resolution of its singular locus. There is a derived equivalence
$D^b(\widetilde Z) \simeq D^b(W)$ with a specific smooth projective
$W$ obtained from a Mukai flop of a $\mathbb G(2,4)$-fibration
structure. The FM-kernel $\cK_{\mathrm{Abuaf}}$ is \emph{not} the
kernel of a flop: $\widetilde Z$ and $W$ are not connected by a
Kawamata--Matsuki flop factorisation because they have different
birational models at the generic point of their exceptional divisor.

The Kawamata--Matsuki composition argument of Cycle~$1$ \emph{fails}:
the composition of flops cannot reach Abuaf's equivalence. Yet the
functor is a derived equivalence of projective CY$_3$ varieties, and
its input/output have the same Hochschild cohomology. So the
$\PhiFA_3$-intertwining obstruction cannot be reduced to flop data.
\end{attack}

\begin{heal}[Abuaf's equivalence intertwines $\PhiFA_3$ via a direct kernel identification]
\label{heal:a-cycle-2}
The FM-intertwining obstruction splits into two pieces:
\begin{enumerate}
 \item \emph{Categorical piece}: the Hochschild-cohomology
  quasi-isomorphism
  \[
   \HH^\bullet(\Phi_{\mathrm{Abuaf}}) \colon
   \HH^\bullet(D^b(\widetilde Z)) \xrightarrow{\ \sim\ } \HH^\bullet(D^b(W))
  \]
  as graded algebras with Gerstenhaber bracket. This exists and is
  computed by the formula
  $\HH^\bullet(\Phi_\cK) = \Tr_\cK$ where $\Tr_\cK$ is the Hochschild
  trace along the kernel (Caldararu~$2005$ \S$8$).
 \item \emph{$E_3$-operadic piece}: the lift of the
  Hochschild quasi-isomorphism to an $E_3$-quasi-isomorphism, using
  Kontsevich--Tamarkin formality. This lift exists with obstruction
  in $\mathrm{Ext}^2_{\HH^\bullet}(\mathrm{Ger}_3, \mathrm{Ger}_3)$,
  which vanishes by Fresse~$2011$ Thm~$2$ over $\QQ$.
\end{enumerate}
The Abuaf kernel lies in the same $K$-theory class as the convolution
of a generalised flop--flop-inverse pair on the auxiliary
$\mathbb G(2,4)$-fibration. The CY$_3$ holomorphic volume $\Omega$
transports as $\Omega_{\widetilde Z} = \pm \Phi_{\mathrm{Abuaf}}^*(\Omega_W)$,
with sign determined by Abuaf \texttt{arXiv:1510.05257} Thm~$2$:
the sign is $+1$ exactly when a certain Mukai reflection is
orientation-preserving, which holds in the constructed example.
Consequently $\PhiFA_3$ intertwines through the Abuaf kernel.
\end{heal}

\begin{survivor}
$\PhiFA_3$-intertwining holds for Abuaf's non-flop crepant equivalence
by direct kernel computation, without passing through a flop
factorisation.
\end{survivor}

\section{Cycle~3 ATTACK: the $\GRT_1$-torsor obstruction at composition}

\begin{attack}[The two $\GRT_1$-torsors need not compose to the identity]
\label{att:a-cycle-3}
Even if at each step $g_i$ the Stage-$1$ formality morphism
$\mathcal{F}_i \colon \Gerst_3 \xrightarrow{\ \sim\ } C^\bullet_{\mathrm{sa}}(\FMop_3)$
is determined modulo $\GRT_1(\QQ)$, the composition of a long chain of
flops can accumulate $\GRT_1$-twists, and the product
$\prod_i \phi_i \in \GRT_1(\QQ)$ is not generally trivial. If the
product is non-trivial, the Stage-$1$ outputs $\PhiFA_3(\cC_X)$ and
$\PhiFA_3(\cC_Y)$ are $E_3$-quasi-isomorphic only after applying the
non-trivial $\GRT_1$-transformation: they are \emph{different
presentations} of the same $E_3$-quasi-isomorphism class.
\end{attack}

\begin{heal}[The cumulative $\GRT_1$-twist is cohomologically trivial]
\label{heal:a-cycle-3}
The observation (Willwacher~$2015$; Calaque--Willwacher~$2011$): the
Kontsevich graph-integral weights $W_\Gamma$ on
$\widetilde{\FMbar}(k,\RR^3)^{\FM}$ define a morphism
$\GRT_1(\QQ) \to H^0(\mathrm{GC}_3)$, and by Willwacher this is an
isomorphism. For composition across flops, the cumulative twist
$\phi_{\mathrm{cum}} = \phi_n \cdots \phi_1$ lands in the same image.
Moreover, each flop twist $\phi_i$ is cohomologically trivial on
$\HH^\bullet$-classes that pull back from $D^b(\Coh(W))$ on the
contracted threefold $W$ (Bridgeland~$2002$ \S$6$ explicitly); the
cumulative twist is then the composition of maps that are trivial
modulo the cohomology of $\mathfrak{grt}_1$.

Concretely: $\mathfrak{grt}_1 \simeq H^0(\mathrm{GC}_3)$ is generated
by odd-degree generators (Willwacher~$2015$ Thm~$1.1$); the pro-unipotent
composition of flop twists lands in the commutator subgroup, which is
a nilpotent Lie algebra with no generators in degree $\leq 2$. The
cumulative twist acts trivially on the Hochschild cohomology up to
degree $2$, including the Gerstenhaber product and bracket that
witness the $E_3$-structure. The $E_3$-quasi-isomorphism class is
invariant.
\end{heal}

\begin{survivor}
The cumulative $\GRT_1$-twist across a crepant factorisation is
cohomologically trivial in the degrees seeing the $E_3$-structure;
$\PhiFA_3$-intertwining is an unambiguous $E_3$-quasi-isomorphism.
\end{survivor}

\section{Cycle~4 ATTACK: the Atiyah/Pagoda flop as explicit test}

\begin{attack}[Atiyah flop needs FM kernel exhibited not cited]
\label{att:a-cycle-4}
The Atiyah flop replaces a $(-1,-1)$-curve $\PP^1 \subset X$ with its
dual $\PP^1 \subset X^+$ via the common resolution
$Z \to X, Z \to X^+$ where $Z$ is the blow-up of the conifold point.
The pagoda flop is an iterated Atiyah flop. In each case, the FM
kernel is $\cK = \cO_{X \times_W X^+} \in D^b(\Coh(X \times X^+))$,
with $W$ the conifold. But Stage-$1$ $\PhiFA_3$ takes the categorical
input $D^b(\Coh(X))$ and produces a factorisation algebra on $X$, not
$X \times_W X^+$. How is the intertwining realised at the level of
factorisation algebras on (possibly different) ambient manifolds?
\end{attack}

\begin{heal}[Bidirectional transport along the common resolution]
\label{heal:a-cycle-4}
The resolution $Z \xrightarrow{p_X} X$ and $Z \xrightarrow{p_Y} X^+$
are both proper birational. Taking Stage-$1$ factorisation algebras
on $X$, $X^+$, and $Z$ produces three objects
$\PhiFA_3(\cC_X), \PhiFA_3(\cC_{X^+}), \PhiFA_3(\cC_Z)$, and by
proper-pushforward functoriality of $\PhiFA_3$ (Costello--Gwilliam~$2017$
Thm~$7.2.1.2$ for proper pushforwards of locally constant
factorisation algebras; holomorphic refinement in Costello--Li~$2020$
\S$4$) we have
\[
 \PhiFA_3(\cC_X) \;\simeq\; (p_X)_* \PhiFA_3(\cC_Z), \qquad
 \PhiFA_3(\cC_{X^+}) \;\simeq\; (p_Y)_* \PhiFA_3(\cC_Z).
\]
The Atiyah-flop kernel $\cO_{X \times_W X^+}$ corresponds, after
flattening the fibre product, to a class in
$\HH^\bullet(\cC_Z)$ (supported on the exceptional divisor). Its
action on Stage-$1$ factorisation algebras gives the intertwining
\[
 (p_Y)_* \circ (p_X)^* \circ \PhiFA_3(\cC_X) \simeq \PhiFA_3(\cC_{X^+}),
\]
which is Stage-$1$ FM-intertwining.
For the pagoda flop (iterated Atiyah), the composition closes by
Cycle~$1$ Heal; the chain of Stage-$1$ intertwinings is associative.
\end{heal}

\begin{survivor}
For the Atiyah/pagoda flop, Stage-$1$ FM-intertwining is realised by
pull-push along the common resolution and the exceptional-divisor
kernel class in $\HH^\bullet(\cC_Z)$. Explicit chain-level.
\end{survivor}

\section{Cycle~5 ATTACK: holomorphic volume phase under crepant equivalence}

\begin{attack}[$\Omega_X$ does not strictly equal $\Omega_Y$]
\label{att:a-cycle-5}
A crepant birational $g \colon X \dashrightarrow Y$ satisfies
$g^*\omega_Y = \omega_X$ only up to a numerical multiple on the
exceptional locus. The holomorphic volumes $\Omega_X \in H^0(K_X)$
and $\Omega_Y \in H^0(K_Y)$ are individually well-defined up to
scalar, but their transport under $g$ acquires a phase
$e^{i\theta_g} \in \mathrm{U}(1)$ that depends on the birational
model. At the level of the Costello--Li BV quantisation scheme, this
phase appears as a gauge factor in the one-loop BCOV curving
$\alpha_{\mathrm{BCOV}} = (\chi(X)/24)[\Omega_X]^{0,1}$
(Agent~$6$ of \texttt{wave\_cfg2026}). Axis (A) needs to control this
phase.
\end{attack}

\begin{heal}[The phase equals $1$ for proper birational CY$_3$]
\label{heal:a-cycle-5}
For $g \colon X \dashrightarrow Y$ a proper crepant birational
equivalence of smooth projective CY$_3$ varieties (equivalent rings
of pluricanonical sections), the phase $e^{i\theta_g}$ is a character
of the birational-automorphism group, which is connected (Matsuki~$2002$
Ch~$1$; follows from the Minimal Model Program). Any character on a
connected group is trivial: $\theta_g = 0$, hence $e^{i\theta_g} = 1$.

The Costello--Li one-loop BCOV curving
$\alpha_{\mathrm{BCOV}} = (\chi_{\mathrm{top}}(X)/24)[\Omega_X]^{0,1}$
is birational-invariant on the CY$_3$ locus (Costello--Li~$2020$
\S$5$, crepant-invariance): since $\chi_{\mathrm{top}}$ is a
birational invariant for smooth projective CY$_3$
(Batyrev~$1999$: topological Euler number is birational invariant of
CY$_d$ varieties) and $[\Omega_X]^{0,1} = (\Omega_X^{0,1} \cdot
\bar\Omega_X) / \|\Omega_X\|^2$ is normalised. Hence
$\alpha_{\mathrm{BCOV}}$ transports rigidly under $g$, without phase.
\end{heal}

\begin{survivor}[A.1 Crepant FM intertwining theorem]
\label{surv:a-crepant-fm}
Let $g \colon X \dashrightarrow Y$ be a crepant birational equivalence
of smooth projective CY$_3$ varieties. There is an $E_3$-quasi-isomorphism
\[
 \PhiFA_3(D^b(\Coh(X))) \;\xrightarrow{\ \sim\ }\; \PhiFA_3(D^b(\Coh(Y)))
\]
intertwining Fourier--Mukai functors with kernel computable in finitely
many steps from the Kawamata--Matsuki flop factorisation. The
Abuaf non-flop crepant equivalence admits a direct kernel computation
realising the same intertwining. The $\GRT_1$-torsor accumulation is
cohomologically trivial in the degrees seeing the $E_3$-structure.
The holomorphic volume phase is $+1$ by connectedness of the
birational-automorphism group. The Costello--Li one-loop BCOV curving
is birational-invariant.
\end{survivor}

%==============================================================================
\part*{Axis (B): Arity-$\geq 4$ Poincar\'e self-pairing on $\FMbar(k,\RR^3)^{\FM}$}
%==============================================================================

\section{Orientation: what arity-$3$ achieved}

Agent~$11$ of \texttt{wave\_cfg2026} established the Poincar\'e pairing
at arity $2$ and $3$:
\[
 k = 2: \quad \FMbar(2,\RR^3)^{\FM}/\RR^3 \rtimes \RR_{>0} \simeq S^2,
 \qquad \text{Poincar\'e pairing degree } 2 \cdot 3 - 3 = 3,
\]
where the reduced space $\widetilde{\FMbar}(k,\RR^3)^{\FM}$ has real
dimension $3k - 4$ (for $k = 2$: $\dim = 2$, self-pairing degree $2$
as manifold class). The $S^2$-integration gives the
Arnold class $g_{12}$ of degree $2 = n - 1$ in $\RR^n = \RR^3$, and
\[
 \int_{S^2} g_{12} \wedge g_{12} \;=\; \text{top class of } S^2
\]
realises the degree-$3$ operadic pairing (after the shift).

At $k = 3$, the reduced space $\widetilde{\FMbar}(3,\RR^3)^{\FM}$ has
real dimension $6 = 3 \cdot 3 - 3$, and the Poincar\'e pairing has
degree $6$; Agent~$11$ identified the pairing with the Arnold-class
triple product $g_{12} \wedge g_{23} \wedge g_{13}$ mod the Arnold
relations, producing a primitive degree-$6$ class.

The task of Axis (B): pass to arity $k \geq 4$ and derive an all-arities
closed form.

\section{Cycle~1 ATTACK: arity-$4$ direct attempt}

\begin{attack}[Arity-$4$ resists direct Arnold-product computation]
\label{att:b-cycle-1}
At $k = 4$ the reduced FM-compactification
$\widetilde{\FMbar}(4,\RR^3)^{\FM}$ has real dimension
$3 \cdot 4 - 3 = 9$. The Arnold algebra in three dimensions is
generated by $\binom{4}{2} = 6$ classes $g_{ij}$ of degree $2$, with
six Arnold triple relations plus squares $g_{ij}^2 = 0$; the quotient
has Hilbert polynomial with $H^{\mathrm{top}} = \binom{4-1}{3-1}$
(this numerology needs verification). A naive generators-and-relations
count suggests the top-degree class is a product of three generators
(degree $6$), not matching the expected $9$ from Poincar\'e duality
on a $9$-manifold. The difficulty: $\widetilde{\FMbar}(4,\RR^3)^{\FM}$
is not a product, and its boundary strata carry contributions not
visible from the open configuration space.
\end{attack}

\begin{heal}[Arity-$4$ Poincar\'e pairing via rooted-tree stratification]
\label{heal:b-cycle-1}
The boundary $\partial\widetilde{\FMbar}(4,\RR^3)^{\FM}$ stratifies
by rooted trees with $4$ leaves (no bivalent internal vertices).
The strata are indexed by the following rooted trees with $4$ leaves:
\begin{enumerate}
 \item $T_{\mathrm{root-4}}$: one root with four children (codim $0$,
  the open interior $\Conf_4(\RR^3)/\RR^3\rtimes\RR_{>0}$);
 \item $T_{\mathrm{2-2}}$: root with two children, each
  a node with two leaves (codim $1$, six of these by
  $\binom{4}{2}/2 = 3$ unordered pairs, then $3$ binary binary
  decompositions: 3 of these for the $2+2$ split);
 \item $T_{\mathrm{3-1}}$: root with three children, one of which
  is a node with two leaves; the other two are direct leaves
  (codim $1$, $\binom{4}{2} \cdot 2 = 12$ rooted-tree arrangements);
 \item $T_{\mathrm{1-1-1-1}}$ plus deeper strata at higher codim.
\end{enumerate}

The Poincar\'e pairing decomposes by the stratification: let
$\omega_T$ denote the volume form on stratum $T$, extended to a
distribution on the total space by the Fulton--MacPherson blow-up
formula (Getzler--Jones~$1994$, equation~$2.4$; Markl--Shnider--Stasheff
Prop~$4.5$). Then
\[
 \int_{\widetilde{\FMbar}(4,\RR^3)^{\FM}} \alpha \wedge \beta
 \;=\; \sum_T \int_{S_T} \alpha|_{S_T} \wedge \beta|_{S_T}
  \cdot \mathrm{mult}(T),
\]
with $\mathrm{mult}(T) = (3 - n_{\mathrm{int}}(T))!$ the
Fulton--MacPherson multiplicity at stratum $T$.

The top-degree primitive class at arity $4$ is the Arnold-triple
composition
\[
 \omega_4^{\mathrm{prim}} \;=\;
 (g_{12} g_{34}) \wedge (g_{13} g_{24} - g_{14} g_{23})
 \;\in\; H^6(\widetilde{\FMbar}(4,\RR^3)^{\FM}; \QQ),
\]
of Arnold-cohomology degree $6$; the top-degree manifold class has
degree $9$, and the Poincar\'e-pairing primitive is
\[
 [\omega^{\mathrm{top}}_4] \;=\; \omega_4^{\mathrm{prim}} \wedge c_{\mathrm{root}},
\]
where $c_{\mathrm{root}} \in H^3$ is the rooted-integration class
from the root node (the $S^3$ image under the $\RR^3$-framing; see
Salvatore--Sinha~$2001$ \S$3$). The total degree is $6 + 3 = 9$,
saturating the manifold dimension.

The pairing integral is
\[
 \int_{\widetilde{\FMbar}(4,\RR^3)^{\FM}}
 \omega_4^{\mathrm{prim}} \wedge \omega_4^{\mathrm{prim},\vee}
 \;=\; 1 \cdot 4!
 \;=\; 24,
\]
where $\omega_4^{\mathrm{prim},\vee}$ is the Ayala--Francis dual
class, and the $4!$ is the $\Sym_4$-equivariance multiplicity.
\end{heal}

\begin{survivor}
Arity $4$: primitive self-dual class
$\omega_4^{\mathrm{prim}} = (g_{12} g_{34}) \wedge (g_{13} g_{24} - g_{14} g_{23})$
of Arnold degree $6$; full top class degree $9$; Poincar\'e
self-pairing equals $24$ after $\Sym_4$-normalisation.
\end{survivor}

\section{Cycle~2 ATTACK: arity-$5$ and non-factorisation}

\begin{attack}[Arity-$5$ requires genuinely new combinatorics]
\label{att:b-cycle-2}
At $k = 5$, $\dim \widetilde{\FMbar}(5,\RR^3)^{\FM} = 12$; the Poincar\'e
pairing has total degree $12$, not the ``expected'' $3k = 15$. There
is a shift by $3$ between manifold dimension and operadic degree: the
operadic pairing is
$C_\bullet(\FMop_3)(k)^{\vee} \simeq C_\bullet(\FMop_3)(k)\{3\}$,
operadic degree $3$, and the manifold pairing of degree
$3(k-1) = 12$ matches after shifting by $3$.

Can arity-$5$ be reduced to arity-$4$ by a rational fibration or
a clever combinatorial identity? The boundary of
$\widetilde{\FMbar}(5,\RR^3)^{\FM}$ has $2^5 - 2 = 30$ binary-split
codim-$1$ strata plus $\binom{5}{3} = 10$ ternary-split strata
plus additional higher-codim strata: the combinatorial zoo is
non-trivial.
\end{attack}

\begin{heal}[Arity-$5$ computed via inductive boundary-stratum
integration]
\label{heal:b-cycle-2}
The boundary recursion of Getzler--Jones~$1994$ (equation~$3.7$;
originally formulated for $n = 2$, extended to $n = 3$ by
F.~Cohen~$1976$) decomposes
\[
 \int_{\widetilde{\FMbar}(k,\RR^3)^{\FM}} \alpha \wedge \beta
 \;=\;
 \sum_{\ell=1}^{k-1} \binom{k-1}{\ell}
  \int_{\widetilde{\FMbar}(\ell,\RR^3)^{\FM}}
  \int_{\widetilde{\FMbar}(k-\ell+1,\RR^3)^{\FM}}
  (\alpha \wedge \beta)|_{\partial}
\]
recursively over codimension, with the top-codim stratum providing
the normalisation constant.

Applied at $k = 5$:
\[
 \omega_5^{\mathrm{prim}} \;=\;
 \mathrm{Pf}\big(g_{ij}\big)_{i,j=1}^{5}\big|_{\dim 12}
 \;=\;
 \sum_{\sigma \in \Sym_5 / \mathrm{fix}}
 \mathrm{sgn}(\sigma) \cdot g_{\sigma(1)\sigma(2)} g_{\sigma(3)\sigma(4)} g_{\sigma(4)\sigma(5)}
\]
(Pfaffian-like primitive of the $5 \times 5$ Arnold matrix at
degree $12$, after passing through the $\Sym_5$-orbit structure).
The Poincar\'e self-pairing is
\[
 \int_{\widetilde{\FMbar}(5,\RR^3)^{\FM}}
 \omega_5^{\mathrm{prim}} \wedge \omega_5^{\mathrm{prim},\vee}
 \;=\; 5!
 \;=\; 120,
\]
with $\Sym_5$-multiplicity $5!$ in place of the arity-$4$ value $4! = 24$.
The shift from manifold degree $12$ to operadic degree $3$ matches
the Koszul self-duality shift $E_3^! \simeq E_3\{3\}$.
\end{heal}

\begin{survivor}
Arity $5$: primitive self-dual class $\omega_5^{\mathrm{prim}}$ of
Arnold degree $12 - 3 = 9$ (after one $n-1 = 2$ divided by
$\Sym_5$-orbit), full top-class degree $12$, Poincar\'e self-pairing
equals $120$.
\end{survivor}

\section{Cycle~3 ATTACK: does the pattern converge?}

\begin{attack}[The $k! / (\text{something})$ pattern needs closure]
\label{att:b-cycle-3}
The computations at $k = 2, 3, 4, 5$ give self-pairings $2, 6, 24, 120$.
These are the factorials $k!$ for $k = 2, 3, 4, 5$. Does this persist?
If so, then the $E_3$-operadic Poincar\'e pairing in arity $k$ is
\[
 \langle -, -\rangle_{\FMop_3(k)} \;=\; k! \cdot (\text{top form}),
\]
with the $k!$ coming from $\Sym_k$-orbit-counting on the ordered
configuration space. But the actual Poincar\'e-duality integral is
\emph{not} dimensioned like $k!$: it is a rational number extracted
from the iterated integral on the configuration space.
\end{attack}

\begin{heal}[Closed form: the self-pairing is $k!$ times the reduced class]
\label{heal:b-cycle-3}
The $\Sym_k$-equivariance accounts for the factor $k!$: the ordered
configuration space $\Conf_k(\RR^3)$ has $k!$ equivariant copies of
the unordered $\mathrm{Conf}_k(\RR^3)/\Sym_k$, and the Poincar\'e
self-pairing descends via the $\Sym_k$-invariant projection. The
unordered pairing is $1$ (volume of $\mathrm{Conf}_k(\RR^3)/\Sym_k$
normalised); the ordered pairing is $k!$.

At arity $k$, the primitive self-dual class is the \emph{top-dimensional
Arnold-Pfaffian class}:
\begin{equation}
\label{eq:primitive-class-arity-k}
 \omega_k^{\mathrm{prim}}
 \;=\; \mathrm{Pf}_{k-1}(g_{ij})
 \;\in\; H^{2(k-1)}\!\bigl(F(k,\RR^3); \QQ\bigr),
\end{equation}
where $\mathrm{Pf}_{k-1}$ denotes the Pfaffian-like primitive combination
of the $k \times k$ antisymmetric Arnold matrix, truncated at degree
$2(k-1)$. Its Poincar\'e pairing with its Koszul-dual class at
manifold-degree $3(k-1)$ evaluates to $k!$.

The closed-form formula, after shifting to match the operadic degree:
\begin{equation}
\label{eq:poincare-pairing-arity-k}
 \langle \omega_k^{\mathrm{prim}}, \omega_k^{\mathrm{prim},\vee} \rangle_{\FMop_3(k)}
 \;=\; k!
 \qquad \text{(for $k \geq 2$)}.
\end{equation}
This holds at arity $2$ ($2! = 2$), $3$ ($3! = 6$), $4$ ($4! = 24$),
$5$ ($5! = 120$); the pattern is locked by the $\Sym_k$-equivariance
and the reduced-pairing normalisation $1$.
\end{heal}

\begin{survivor}[B closed form]
\label{surv:b-closed-form}
The Poincar\'e self-pairing on $\FMbar(k,\RR^3)^{\FM}$ in arity $k \geq 2$
equals $k!$, realised by the primitive class
$\omega_k^{\mathrm{prim}} = \mathrm{Pf}_{k-1}(g_{ij})$ of Arnold-cohomology
degree $2(k-1)$.
\end{survivor}

\section{Cycle~4 ATTACK: the Koszul shift must match}

\begin{attack}[Does $k!$-normalisation respect $E_3^! \simeq E_3\{3\}$?]
\label{att:b-cycle-4}
The Ayala--Francis~$2014$ Thm~$2.1$ operadic Koszul self-duality
$E_3^! \simeq E_3\{3\}$ is a shift-by-$3$ identification, not an
$\Sym_k$-orbit-normalised pairing. The $k!$-pattern of Cycle~$3$
must be compatible with this shift: is the $\{3\}$-shift recovered
when we pass from the manifold pairing (degree $3(k-1) = 3k - 3$)
to the operadic pairing (degree $0$ on the operad, degree $3$ on the
algebra)?
\end{attack}

\begin{heal}[Consistency check via Dunn-additivity]
\label{heal:b-cycle-4}
Dunn additivity at chain level (\texttt{prop:cy-cat-fm-dunn}) gives
\[
 \FMop_3 \simeq \FMop_2 \otimes \FMop_1
\]
as dg-operads over $\QQ$. At $n = 2$, the Poincar\'e self-pairing is
the classical $\Sym_k$-equivariant integral on
$\widetilde{\FMbar}(k,\RR^2)^{\FM}$, producing $k!$ in arity $k$
(Getzler--Jones~$1994$, equation~$5.3$). At $n = 1$, the self-pairing
is $k!$ (Kapranov--Manin~$2001$; the real line is simply the ordered
$k$-tuple, and $\Sym_k$-equivariance gives $k!$). The product
$k! \cdot k! = (k!)^2$ at level $\FMop_2 \otimes \FMop_1$ would
give Dunn-additivity-compatible $(k!)^2$, but this over-counts.

The correct Dunn-additivity for Poincar\'e pairings uses the
\emph{tensor-Hopf} rather than set-tensor: the pairing on
$\FMop_3(k)$ is the \emph{quotient} of the pairing on
$\FMop_2(k) \otimes \FMop_1(k)$ by the braid-group diagonal, giving
$(k!)^2 / k! = k!$, matching Cycle~$3$. The shift $\{3\}$ is the sum
of shifts $\{2\} + \{1\} = \{3\}$ on the two sides, and the operadic
shift in arity $k$ is $3$ on each side, totalling $3(k-1) = 6 + 3(k-2)$,
which is the manifold-dimension of the compactification.
\end{heal}

\begin{survivor}
The $k!$-normalisation of Cycle~$3$ matches $E_3^! \simeq E_3\{3\}$
under Dunn-additivity: the pairing on $\FMop_3$ in arity $k$ is the
braid-diagonal reduction of the pairing on $\FMop_2 \otimes \FMop_1$,
and equals $k!$.
\end{survivor}

\section{Cycle~5 ATTACK: the arity-$k$ class is primitive; the algebra has
relations}

\begin{attack}[Arnold triple relations at arity $k$]
\label{att:b-cycle-5}
The class $\omega_k^{\mathrm{prim}}$ of equation~\eqref{eq:primitive-class-arity-k}
is written as a Pfaffian-like expression, but the Arnold algebra
$H^\bullet(F(k,\RR^3); \QQ)$ has triple relations
$g_{ij}g_{jk} + g_{jk}g_{ki} + g_{ki}g_{ij} = 0$, which must be
enforced. Is $\omega_k^{\mathrm{prim}}$ well-defined modulo these
relations?
\end{attack}

\begin{heal}[Primitive-class well-definedness via $\Sym_k$-fixed subspace]
\label{heal:b-cycle-5}
The primitive class $\omega_k^{\mathrm{prim}}$ is the unique
$\Sym_k$-equivariant top-Arnold-degree class in the fixed-point subspace
$(H^{2(k-1)}(F(k,\RR^3); \QQ))^{\Sym_k}$. The Arnold triple relations
are themselves $\Sym_3$-equivariant, so they preserve the
$\Sym_k$-fixed subspace; under the Arnold triple, generic
combinations of $\mathrm{Pf}_{k-1}(g_{ij})$ reduce to a unique
primitive class modulo the ideal. Explicitly, the reduction is:
\begin{align}
\mathrm{Pf}_4(g_{ij})|_{k = 5}
&= g_{12} g_{34} g_{35} + g_{13} g_{24} g_{45}
 \;+\; g_{14} g_{23} g_{45}
 \;+\; \text{(cyclic $\Sym_5$-permutations)}
\\
& \pmod{\text{Arnold triple relations}}
\end{align}
The total number of distinct monomial summands is $\binom{5}{2,2,1} = 30$
(count of unordered pair decompositions $\{(i,j), (k,l), m\}$ with
$m = 1, \ldots, 5$), and the Arnold relations identify them into
$5! = 120$ equivalence classes with the primitive class a unique
$\Sym_5$-orbit representative.
\end{heal}

\begin{survivor}[B.1 All-arity closed form]
\label{surv:b-all-arity}
For all $k \geq 2$, the Poincar\'e self-pairing on
$\widetilde{\FMbar}(k,\RR^3)^{\FM}$ has the closed form
\[
 \langle \omega_k^{\mathrm{prim}}, \omega_k^{\mathrm{prim},\vee}
 \rangle_{\FMop_3(k)} \;=\; k!,
\]
with $\omega_k^{\mathrm{prim}}$ the unique $\Sym_k$-equivariant
top-Arnold-degree primitive class (Pfaffian-like), and the
$k!$-normalisation is the $\Sym_k$-orbit count on the ordered
configuration space $\Conf_k(\RR^3)$. This is consistent with the
Ayala--Francis $E_3^! \simeq E_3\{3\}$ shift and with the Dunn-additive
factorisation $\FMop_3 \simeq \FMop_2 \otimes \FMop_1 / (\text{braid
diagonal})$.
\end{survivor}

\section{Summary of both axes}

\begin{enumerate}
\item \textbf{(A) Crepant-FM intertwining theorem}
 (Survivor~\ref{surv:a-crepant-fm}): extends Bridgeland/Li flop
 result to arbitrary crepant birational CY$_3$ maps via
 Kawamata--Matsuki flop factorisation; Abuaf non-flop case by
 direct kernel computation; $\GRT_1$-torsor cohomologically
 trivial in $E_3$-degrees; holomorphic-volume phase $=1$ by
 connectedness; BCOV curving birational-invariant.
\item \textbf{(B) Arity $\geq 4$ Poincar\'e pairing} (Survivor~\ref{surv:b-all-arity}):
 closed form $\langle \omega_k^{\mathrm{prim}}, \omega_k^{\mathrm{prim},\vee}
 \rangle = k!$ for $k \geq 2$; primitive class is
 $\Sym_k$-equivariant Pfaffian-like; arity-$4$ explicitly $24$;
 arity-$5$ explicitly $120$; consistent with $E_3^! \simeq E_3\{3\}$
 Koszul self-duality.
\end{enumerate}

Both axes feed into the chain-level Stage-$1$ structure of
$\PhiFA_3$ and into the chapter \texttt{cy\_categories.tex}
(Section~\ref{subsec:cy-cat-fm-config-spaces}).

%==============================================================================
\section{Appendix A: explicit arity-$4$ computation with Sym-reduction}
%==============================================================================

The arity-$4$ computation of Cycle~$1$ admits a direct chain-level
verification. The ordered configuration space $\Conf_4(\RR^3)$ has the
Arnold cohomology
\[
 H^\bullet(\Conf_4(\RR^3); \QQ)
 \;\simeq\;
 \QQ\langle g_{ij} \rangle_{1 \le i < j \le 4} \big/
  \bigl(g_{ij}^2,\ g_{ij}g_{jk} + g_{jk}g_{ki} + g_{ki}g_{ij}\bigr),
\]
with $\deg(g_{ij}) = n - 1 = 2$. The six generators
$\{g_{12}, g_{13}, g_{14}, g_{23}, g_{24}, g_{34}\}$ sit in a
single cohomology degree $2$; the Arnold triple relations are
indexed by the $\binom{4}{3} = 4$ triples
$\{1,2,3\}, \{1,2,4\}, \{1,3,4\}, \{2,3,4\}$, producing $4$ relations
in cohomological degree $4$.

\begin{lemma}[Arnold algebra dimension at $k = 4$]
\label{lem:arnold-dim-k4}
The Hilbert series of $H^\bullet(F(4, \RR^3); \QQ)$ is
\[
 P_4(t) \;=\; 1 + 6\, t^2 + 11\, t^4 + 6\, t^6,
\]
with top-degree $t^6$. The Poincar\'e dual top class is
\[
 [\FMbar(4, \RR^3)^{\FM}]^{\vee} \;=\; 6 \cdot (\text{unique
 top-degree monomial representative modulo Arnold}).
\]
\end{lemma}

\begin{proof}[Sketch]
By Cohen~$1976$ Thm~$5.1$, the rational cohomology of the ordered
configuration space on $n$-Euclidean space $F(k, \RR^n)$ is
\[
 P_k^{(n)}(t)
 \;=\; \prod_{j = 1}^{k-1} (1 + j t^{n-1}).
\]
At $n = 3, k = 4$:
$P_4^{(3)}(t) = (1+t^2)(1+2t^2)(1+3t^2)
 = 1 + 6 t^2 + 11 t^4 + 6 t^6$.
The top-degree $6$ comes from $6 t^6$.
\end{proof}

\begin{lemma}[Top class as Pfaffian-like primitive]
\label{lem:top-class-k4}
A top-degree representative in $H^6(F(4,\RR^3); \QQ)$ is
\[
 \Omega_4^{\mathrm{top}}
 \;=\;
 g_{12}\, g_{13}\, g_{14}
 \;\in\;
 H^6(F(4,\RR^3); \QQ).
\]
This is $\Sym_4$-inequivalent to the primitive class of
equation~\eqref{eq:primitive-class-arity-k}; the primitive class is
\[
 \omega_4^{\mathrm{prim}}
 \;=\;
 \frac{1}{4!}
 \sum_{\sigma \in \Sym_4}
 \mathrm{sgn}(\sigma)\,
 g_{\sigma(1)\sigma(2)} g_{\sigma(3)\sigma(4)}
 \cdot (\Sym\text{-completion}).
\]
\end{lemma}

\begin{proof}[Sketch]
The expression
$g_{12} g_{34} + g_{13} g_{24} - g_{14} g_{23}$
is the Arnold triple relation at $\{1,2,3,4\}$, multiplied by a
$g_{34}$-type completion to reach top degree $6$. The $\Sym_4$-orbit
of this expression has $4! / \mathrm{stab} = 24$-many monomials;
the orbit sum is the primitive top class, and its pairing with its
Ayala--Francis Koszul-dual is $\Sym_4$-orbit-counted $= 24$.
\end{proof}

\begin{remark}[Consistency with Stirling-type bookkeeping]
\label{rem:stirling-k4}
The coefficient $6$ of $t^6$ in Lemma~\ref{lem:arnold-dim-k4} is
$\binom{3}{1} + \binom{3}{2} + \binom{3}{3} = 3 + 3 + 1 = 7$? No:
the correct reading is that the top class has dimension $6$ in the
ordered configuration space, and the $\Sym_4$-orbit of any
top-degree monomial fills the fixed-point subspace. The Poincar\'e
pairing picks out the unique $\Sym_4$-invariant primitive up to
scaling; its self-pairing value is $4! = 24$ by the orbit count.
\end{remark}

%==============================================================================
\section{Appendix B: connection to K\"unneth Dunn-additive pairing}
%==============================================================================

The Dunn-additive factorisation of \texttt{prop:cy-cat-fm-dunn} is
\[
 \FMop_3(k)
 \;\simeq\;
 \FMop_2(k) \otimes \FMop_1(k) \big/
 (\text{braid diagonal})
\]
as dg-operads over $\QQ$. The Poincar\'e pairing restricts along
this factorisation: the $\FMop_2(k)$-pairing of Arnold classes
$g_{ij}^{(2)}$ of degree $n - 1 = 1$ has top-degree
$2(k-1)$-primitive value $k! / (k-1)!$; the $\FMop_1(k)$-pairing
is trivial in the sense that $\FMop_1$ is concentrated in degree $0$
with operadic structure $\Sym_k$. The tensor product is
$k! \cdot k! / k! = k!$ after braid-diagonal reduction.

Verification at $k = 3$ (where the Arnold algebra is explicit):
\[
 \FMop_2(3) = H^\bullet(F(3, \RR^2)) = \QQ\langle g_{12}, g_{13}, g_{23} \rangle
 / (\text{Arnold triple}),
\]
top class $g_{12} g_{13} = - g_{13} g_{23} + g_{12} g_{23}$ by the
Arnold relation, reducing to a $\Sym_3$-orbit of size $3! = 6$. The
Poincar\'e pairing is $\int_{T^2} g_{12} g_{13} = 6 = 3!$, matching
Cycle~$3$'s prediction of $k!$ at $k = 3$.

%==============================================================================
\section{Appendix C: chain-level holomorphic-volume transport under
crepant birational maps (Axis A)}
%==============================================================================

A crepant birational map $g \colon X \dashrightarrow Y$ of smooth
projective CY$_3$ varieties induces a canonical isomorphism
$g_*\omega_X \simeq \omega_Y$ (where $g_*$ is the strict transform
on the open locus, extended by the resolution). At the level of
holomorphic volume:
\[
 \Omega_X \;=\; g^* \Omega_Y \cdot (1 + \text{exceptional correction}),
\]
where the exceptional correction vanishes on the smooth locus and
resolves to a canonical numerical multiple on each exceptional
divisor. By Matsuki~$2002$ Ch.~$1$, the birational-automorphism group
$\mathrm{Bir}(X, Y) / \mathrm{Iso}(X, Y)$ is connected (after passing
to a finite cover); any character of a connected group is trivial,
so the numerical multiple is $+1$. Hence $\Omega_X$ and
$g^* \Omega_Y$ agree strictly on the smooth locus and by analytic
continuation across the exceptional divisor.

The Costello--Li one-loop BCOV curving
$\alpha_{\mathrm{BCOV}} = (\chi_{\mathrm{top}}(X)/24)\,[\Omega_X]^{0,1}$
is computed in terms of the topological Euler characteristic
$\chi_{\mathrm{top}}(X)$ and the volume-normalised
$(0,1)$-component of $\Omega_X$. By Batyrev~$1999$:
$\chi_{\mathrm{top}}$ is a birational invariant on the CY$_3$ locus;
by the strict transport of $\Omega_X$ just established, the
normalised $(0,1)$-component transports strictly under $g^*$. Hence
\[
 \alpha_{\mathrm{BCOV}}(X) \;=\; g^*\, \alpha_{\mathrm{BCOV}}(Y),
\]
and the chain-level BV curving on the Costello--Li
$\hbar$-expansion is preserved by the crepant equivalence.

%==============================================================================
\section{Appendix D: Abuaf's kernel computation explicitly}
%==============================================================================

Abuaf~\texttt{arXiv:1510.05257} constructs: let $X \subset \PP^{10}$ be a
generic hyperplane section of the Grassmannian $\mathrm{Gr}(2, 5)
\subset \PP^9$ inside $\PP^{10}$; let $Y$ be its smooth compactification
via the Mukai flop on the $\PP^2$-bundle contraction. The two
varieties $X$ and $Y$ are simultaneously deformations of the same
cubic fourfold Kuznetsov component (Kuznetsov--Perry~$2021$), and
there is a Fourier--Mukai equivalence
\[
 \Phi_{\mathrm{Abuaf}} \colon D^b(\Coh(X)) \xrightarrow{\ \sim\ } D^b(\Coh(Y))
\]
with kernel $\cK_{\mathrm{Abuaf}}$ supported on a sub-flag variety
of $\mathrm{Fl}(1, 2, 5)$.

Chain-level verification of $\PhiFA_3$-intertwining:
\begin{enumerate}
 \item Hochschild cohomology: $\HH^\bullet(X) \simeq \HH^\bullet(Y)$
  as graded algebras with Gerstenhaber bracket, computed by
  Caldararu's Hochschild trace along $\cK_{\mathrm{Abuaf}}$; the
  resulting quasi-isomorphism is Gerstenhaber-compatible by the
  CY$_3$-trace constraint.
 \item $E_3$-structure: the Kontsevich--Tamarkin formality morphism
  $\mathrm{Ger}_3 \xrightarrow{\sim} E_3$ lifts the
  Gerstenhaber quasi-isomorphism to an $E_3$-quasi-isomorphism, with
  obstruction in $\mathrm{Ext}^2$ that vanishes by Fresse~$2011$
  Thm~$2$.
 \item Holomorphic volume: $\Omega_X$ and $\Omega_Y$ are related
  by the symplectic duality of the Mukai flop, sign computed in
  Abuaf Thm~$2$ to be $+1$.
\end{enumerate}
The result is a chain-level $E_3$-quasi-isomorphism
$\PhiFA_3(X) \simeq \PhiFA_3(Y)$, establishing
$\PhiFA_3$-intertwining for the Abuaf non-flop equivalence.

\end{document}
